رغم ارتفاع مستوى عدم اليقين، استمر النمو الاقتصادي العالمي عند 3.3\% بزيادة قدرها 0.2 نقطة مقارنة بتوقعات أكتوبر الماضي. ويُعزى هذا الأداء إلى عدة عوامل، من بينها تراجع حدة التوترات التجارية عقب الهدنة التجارية بين الصين والولايات المتحدة، إلى جانب الحوافز المالية والنقدية، واستمرار الأوضاع المالية التيسيرية. كما ساهمت الطفرة الاستثمارية في قطاع تكنولوجيا المعلومات، ولا سيما في مجال الذكاء الاصطناعي، تحديدا في الولايات المتحدة والصين في دعم هذا النمو. بالإضافة إلى ارتفاع التجارة العالمية في الصادرات ذات الصلة بالتكنولوجيا. ومع ذلك لا تزال معدلات نمو التجارة العالمية متواضعة؛ حيث من المتوقع أن تسجل نحو 2.6\% في عام 2026 وهو مستوى أقل من نظيرتها المسجل في عام 2025.\textsuperscript{4}
